% Рисунок из тетради: «рваный» график u и скользящее под ним ядро.
\begin{figure}[htbp]
\centering
\begin{tikzpicture}
\begin{axis}[
  width=11cm, height=5cm,
  axis lines=middle,
  xmin=-0.2, xmax=6.2,
  ymin=-0.3, ymax=3.4,
  xtick=\empty, ytick=\empty,
  clip=false,
]
  % «рваная» функция u
  \addplot[thick, blue!60!black] coordinates {
    (0.1,0.4) (0.35,2.1) (0.6,0.6) (0.85,2.6) (1.1,1.0) (1.35,1.9)
    (1.6,0.7) (1.9,1.5) (2.2,0.5) (2.5,1.2) (2.8,0.6) (3.1,1.4)
    (3.4,0.8) (3.5,3.0) (3.6,0.9) (3.9,1.1) (4.2,0.5)
    (4.5,3.1) (4.55,0.6) (4.9,2.9) (5.0,0.4) (5.4,1.3) (5.8,0.5)
  };
  \node[blue!60!black, anchor=south] at (axis cs:0.9,2.7) {$u$};

  % ядро под графиком
  \addplot[thick, red!70!black, samples=100, domain=2.1:3.3] {0.9*bump(x-2.7,0.6)};
  \filldraw (axis cs:2.7,0) circle (1.3pt);
  \node[red!70!black, anchor=west] at (axis cs:3.35,0.9) {\footnotesize $y\mapsto\omega_\rho(x-y)$};
  \draw[dashed, gray] (axis cs:2.7,0) -- (axis cs:2.7,1.35);
  \node[anchor=north] at (axis cs:2.7,-0.02) {\footnotesize $x$};
  \draw[-{Latex[length=2mm]}, gray] (axis cs:3.4,0.25) -- (axis cs:4.1,0.25);
  \node[gray, anchor=north, font=\scriptsize] at (axis cs:3.75,0.22) {скольжение по $x$};
\end{axis}
\end{tikzpicture}
\caption{Значение $u_\rho(x)$ — взвешенное среднее $u$ под «шапочкой» радиуса $\rho$
с центром в $x$. Зубцы усредняются, и получается гладкая кривая.}
\label{fig:smoothing}
\end{figure}
